%% annex/p08/modal-logics.tex --- Part 8.2 proof material

\begin{annexitem}{annex:p8.modal.1}{thm:p08-modal-four-system-adequacy}{Canonical-model adequacy for K, T, S4, and S5}
\sketchnote{The Lindenbaum extension lemma for arbitrary premise sets, the
modal diamond/existence lemma, every Boolean and modal case of the truth lemma,
and the full T-, 4-, and 5-to-frame-property verifications are omitted.  The
cited sources supply these results for the exact theorem-only necessitation
discipline and local consequence fixed in the body.}

For soundness, classical tautologies and modus ponens preserve truth at the
chosen world.  The distribution schema is valid on every Kripke frame.
Theorem-only necessitation is sound because a no-premise theorem is valid at
every world before it is boxed; this step would not license necessitation of a
formula depending on an undischarged premise.  The T, 4, and 5 schemas are
valid, respectively, on reflexive, transitive, and Euclidean frames.  In the
S5 row, reflexivity and Euclideanness make the relation an equivalence relation.

For completeness, suppose
\(\Gamma\notderives_{\mathfrak L}A\).  Finitarity of derivations makes
\(\Gamma\cup\{\neg A\}\) \(\mathfrak L\)-consistent.  By the omitted
Lindenbaum lemma, extend it to a maximal \(\mathfrak L\)-consistent set \(w\).
Let the worlds be all maximal \(\mathfrak L\)-consistent sets, set
\[
 uR_{\mathfrak L}v
 \quad\text{ iff }\quad
 \forall B\,(\Box B\in u\to B\in v),
\]
and make an atom true at a world exactly when it belongs to that set.

The omitted existence lemma says that if \(\Diamond B\in u\), then some
\(v\) satisfies \(uR_{\mathfrak L}v\) and \(B\in v\).  Using that lemma in the
diamond direction, structural induction yields the omitted truth lemma
\[
 \mathcal M_{\mathfrak L},u\entails B
 \quad\text{ iff }\quad
 B\in u.
\]
Hence the distinguished world \(w\) satisfies every member of \(\Gamma\) and
does not satisfy \(A\), giving the local countermodel.

It remains to keep the canonical frame inside the theorem's selected class.
For T, if \(\Box B\in u\), the T instance gives \(B\in u\), so
\(uR_{\mathsf T}u\).  For 4, if \(uRv\), \(vRz\), and \(\Box B\in u\), the
4 instance places \(\Box\Box B\) in \(u\), hence \(\Box B\in v\) and
\(B\in z\); thus the canonical relation is transitive.  The omitted 5
verification gives Euclideanness, and together with the T verification gives
an equivalence relation for S5.  K imposes no additional frame condition.
These checks finish the sketched reduction to the omitted canonical lemmas
\cite{BlackburnDeRijkeVenema2001,Chellas1980}.
\end{annexitem}
